\section{Generalizations of the Coherent--Constructible Correspondence}

The previous sections focused on the coherent--constructible correspondence for
smooth toric varieties. In that setting, the input is a fan
\(\Sigma\subset N_{\mathbb R}\), the coherent side is modeled by
\(\mathrm{Perf}_T(X_\Sigma)\), and the constructible side is the category of
constructible sheaves on \(M_{\mathbb R}\) whose singular support is contained
in the FLTZ skeleton \(\Lambda_\Sigma\). The key mechanism is the correspondence
between coherent theta objects and constructible theta sheaves:
\[
\Theta'(\sigma,\chi)
\longleftrightarrow
\Theta(\sigma,\chi).
\]
The Hom-complex comparison is then controlled by inclusions of polyhedra
\[
\chi+\sigma^\vee\subset \psi+\tau^\vee.
\]

The purpose of this section is to indicate several directions in which this
picture can be generalized. We will not give complete proofs of these
generalizations. Instead, we explain how the basic objects of the toric CCC are
modified, and how the same organizing principle remains visible.

There are three generalizations that are especially relevant for the present
discussion. The first replaces toric varieties by toric Deligne--Mumford stacks.
In this case, the fan is replaced by a stacky fan. The constructible side is
still controlled by a conic subset of \(T^*M_{\mathbb R}\), but the lattice data
are modified by the stacky structure. This changes the indexing of theta
objects and records finite stabilizer information.

The second generalization concerns toric fibrations. Instead of considering a
single toric variety over a point, one considers a family whose fibers are
toric varieties with fixed fan. The coherent side becomes a category of sheaves
on a relative toric space over a base, while the constructible side is no
longer just a single sheaf category on \(M_{\mathbb R}\). It is naturally
described in terms of a sheaf, or local system, of categories over the base.

The third generalization concerns Fourier--Mukai transforms. In the usual
coherent setting, Fourier--Mukai kernels define functors between derived
categories of coherent sheaves. Through the CCC, such functors can be translated
into operations on constructible sheaves. This gives a constructible
interpretation of certain derived equivalences and birational transformations
in toric geometry.

Thus, the guiding question of this chapter is not whether the formula for the
CCC remains exactly the same. Rather, it is how each part of the correspondence
changes:
\[
\text{fan}
\quad\rightsquigarrow\quad
\text{modified combinatorial input},
\]
\[
\mathrm{Perf}_T(X_\Sigma)
\quad\rightsquigarrow\quad
\text{coherent category in the new setting},
\]
\[
Sh_{\Lambda_\Sigma}(M_{\mathbb R})
\quad\rightsquigarrow\quad
\text{constructible or microlocal category with modified support condition}.
\]
In each case, the same basic philosophy remains: algebraic data on the coherent
side are translated into constructible sheaves with prescribed microlocal
behavior.


\subsection{Stacky Fans and Toric Deligne--Mumford Stacks}

The first generalization replaces toric varieties by toric Deligne--Mumford
stacks. The main change is not the form of the correspondence, but the
combinatorial input. Instead of an ordinary fan \(\Sigma\subset N_{\mathbb R}\),
one uses a stacky fan
\[
\boldsymbol\Sigma=(N,\Sigma,\beta).
\]
Here \(\Sigma\) still records the real cone directions, while \(\beta\) records
the stacky lattice data along the rays. Thus the constructible side is still
controlled by a conic subset of \(T^*M_{\mathbb R}\), but the indexing of theta
objects is modified by the stacky structure.

The basic mechanism remains the same. Coherent theta objects are matched with
constructible theta sheaves,
\[
\Theta'(\sigma,\chi)\longleftrightarrow \Theta(\sigma,\chi),
\]
and Hom-complexes are still controlled by inclusions of polyhedra. What changes
is the meaning of the character \(\chi\). In the stacky case, the local lattice
attached to a cone may be different from the ordinary lattice coming from the
coarse toric variety.

We illustrate this by the weighted projective line \(\mathbb P(1,2)\). Its
coarse moduli space is \(\mathbb P^1\), but it has a stacky point of order
\(2\). We describe it by the following stacky fan. Let \(N=\mathbb Z\), and let
\(\Sigma\) be the complete fan in \(N_{\mathbb R}=\mathbb R\), with rays
\[
\rho_+=\mathbb R_{\ge 0},
\qquad
\rho_-=\mathbb R_{\le 0}.
\]
The stacky vectors are chosen to be
\[
b_+=1,\qquad b_-=-2.
\]
Equivalently, the map \(\beta:\mathbb Z^2\to N\) sends the two standard basis
vectors to \(1\) and \(-2\). The nonprimitive vector \(-2\) records the stacky
structure at one torus fixed point.

We now compute the corresponding FLTZ skeleton. The zero cone contributes the
zero section
\[
\mathbb R\times\{0\}\subset T^*\mathbb R.
\]
For the ray \(\rho_+\), the local lattice is generated by \(b_+=1\), so the
corresponding character lattice is the ordinary lattice \(\mathbb Z\). Hence
\(\rho_+\) contributes negative conormal rays over the integral points:
\[
\bigcup_{a\in\mathbb Z}\{a\}\times\mathbb R_{\le 0}.
\]
For the ray \(\rho_-\), the local lattice is generated by \(b_-=-2\). Its dual
lattice is naturally identified with \(\frac{1}{2}\mathbb Z\subset\mathbb R\).
Therefore \(\rho_-\) contributes positive conormal rays over the half-integral
points:
\[
\bigcup_{b\in \frac{1}{2}\mathbb Z}\{b\}\times\mathbb R_{\ge 0}.
\]
Thus
\[
\Lambda_{\boldsymbol\Sigma}
=
\mathbb R\times\{0\}
\cup
\bigcup_{a\in\mathbb Z}\{a\}\times\mathbb R_{\le 0}
\cup
\bigcup_{b\in \frac{1}{2}\mathbb Z}\{b\}\times\mathbb R_{\ge 0}.
\]
Compared with the ordinary \(\mathbb P^1\), the new feature is the appearance
of half-integral base points in one conormal direction. This is the
constructible trace of the stacky stabilizer of order \(2\).

The theta sheaves are described in the same way as before, but with the
modified indexing set. For \(a\in\mathbb Z\), the ray \(\rho_+\) gives
\[
P_{\rho_+,a}=[a,\infty),
\qquad
\Theta(\rho_+,a)
=
(i_{[a,\infty)})_*\omega_{[a,\infty)}.
\]
For \(b\in \frac{1}{2}\mathbb Z\), the ray \(\rho_-\) gives
\[
P_{\rho_-,b}=(-\infty,b],
\qquad
\Theta(\rho_-,b)
=
(i_{(-\infty,b]})_*\omega_{(-\infty,b]}.
\]
Thus the stacky structure does not change the fact that theta sheaves are
costandard sheaves on half-lines, but it changes where their endpoints are
allowed to lie.

The Hom-complex calculation has the same form as in the ordinary toric case.
For \(a,a'\in\mathbb Z\),
\[
\mathrm{RHom}(\Theta(\rho_+,a),\Theta(\rho_+,a'))
\simeq
\begin{cases}
k, & a\ge a',\\
0, & a<a',
\end{cases}
\]
because \([a,\infty)\subset [a',\infty)\) if and only if \(a\ge a'\). For
\(b,b'\in \frac{1}{2}\mathbb Z\),
\[
\mathrm{RHom}(\Theta(\rho_-,b),\Theta(\rho_-,b'))
\simeq
\begin{cases}
k, & b\le b',\\
0, & b>b',
\end{cases}
\]
because \((-\infty,b]\subset(-\infty,b']\) if and only if \(b\le b'\).

This example shows the basic point of the stacky generalization. The formula
for the Hom-complex is unchanged; it is still governed by inclusion of
polyhedra. What changes is the lattice of allowed translations. In the ordinary
\(\mathbb P^1\) case, both families of endpoints lie in \(\mathbb Z\). For the
weighted projective line \(\mathbb P(1,2)\), one family still lies in
\(\mathbb Z\), while the other lies in \(\frac{1}{2}\mathbb Z\). The fractional
positions in the FLTZ skeleton encode the stacky structure.

\subsection{Toric Fibrations}


The second generalization replaces a single toric variety by a toric fibration.
The guiding idea is that the fan still controls the toric directions, while
the base of the fibration contributes coefficients to the Hom-complexes.

Let \(B\) be a smooth variety, let \(T\) be an algebraic torus, and let
\(P\to B\) be a principal \(T\)-bundle. Given a fan \(\Sigma\) for a toric
variety \(X_\Sigma\), one can form the associated toric fibration
\[
X_{\Sigma,P}:=P\times_T X_\Sigma \longrightarrow B.
\]
The fiber over each point of \(B\) is isomorphic to \(X_\Sigma\). Thus the
toric geometry is constant along the base, but the gluing over \(B\) may be
twisted by the principal \(T\)-bundle.

The constructible side should be viewed as a relative version of the ordinary
constructible category. The FLTZ skeleton is still determined by the fan
\(\Sigma\), so the microlocal directions in \(T^*M_{\mathbb R}\) are the same
as in the toric variety case. What changes is that the constructible sheaves
now carry coefficients coming from the base \(B\). Equivalently, the ordinary
condition
\[
SS(F)\subset \Lambda_\Sigma
\]
is kept, but the values of the sheaves are allowed to remember objects on
\(B\), such as perfect complexes or line bundles associated to the principal
\(T\)-bundle.

The coherent theta objects are also modified in a simple way. For a character
\(\lambda\in M\), the principal \(T\)-bundle \(P\to B\) determines an associated
line bundle
\[
\mathcal L_\lambda:=P\times_T \mathbb C_\lambda
\]
on \(B\), where \(\mathbb C_\lambda\) denotes the one-dimensional
\(T\)-representation of weight \(\lambda\). Thus a monomial of weight
\(\lambda\) no longer contributes only a scalar; it contributes the line bundle
\(\mathcal L_\lambda\) on the base. This is the main new feature of the
fibration case.

We illustrate this in the simplest example. Let \(T=\mathbb G_m\), let
\(P\to B\) be a principal \(\mathbb G_m\)-bundle, and take \(X_\Sigma=\mathbb P^1\).
The associated space
\[
X_{\Sigma,P}=P\times_{\mathbb G_m}\mathbb P^1
\]
is a \(\mathbb P^1\)-bundle over \(B\).

The fan is the complete fan in \(N_{\mathbb R}=\mathbb R\), with rays
\[
\rho_+=\mathbb R_{\ge 0},
\qquad
\rho_-=\mathbb R_{\le 0}.
\]
Therefore the FLTZ skeleton is the same one as for \(\mathbb P^1\):
\[
\Lambda_\Sigma
=
\mathbb R\times\{0\}
\cup
\bigcup_{a\in\mathbb Z}\{a\}\times\mathbb R_{\le 0}
\cup
\bigcup_{a\in\mathbb Z}\{a\}\times\mathbb R_{\ge 0}.
\]
Thus the allowed singular directions are unchanged. The difference is that the
theta sheaves now carry coefficients from \(B\).

For \(a\in\mathbb Z\), the ray \(\rho_+\) gives the half-line
\[
P_{\rho_+,a}=[a,\infty),
\]
and the corresponding constructible theta object is still modeled on
\[
(i_{[a,\infty)})_*\omega_{[a,\infty)}.
\]
Similarly, the ray \(\rho_-\) gives theta objects supported on half-lines of
the form \((-\infty,a]\). Hence the polyhedral part of the construction is the
same as before.

Now consider the Hom-complex. In the ordinary toric case, one has
\[
\mathrm{RHom}(\Theta(\rho_+,a),\Theta(\rho_+,b))
\simeq
\begin{cases}
k, & a\ge b,\\
0, & a<b.
\end{cases}
\]
For the toric fibration, the same inclusion condition remains, but the
nonzero term is replaced by cohomology on the base. With the convention above,
a morphism from weight \(a\) to weight \(b\) is represented by a monomial of
weight \(a-b\). Therefore
\[
\mathrm{RHom}_{X_{\Sigma,P}}
(\Theta'_P(\rho_+,a),\Theta'_P(\rho_+,b))
\simeq
\begin{cases}
\mathrm{R}\Gamma(B,\mathcal L_{a-b}), & a\ge b,\\
0, & a<b.
\end{cases}
\]
The condition \(a\ge b\) is exactly the same condition as
\[
[a,\infty)\subset [b,\infty).
\]
Thus the polyhedral inclusion still determines whether the Hom-complex
vanishes, while the base \(B\) determines the value of the nonzero
Hom-complex.

The same calculation applies to the other ray:
\[
\mathrm{RHom}_{X_{\Sigma,P}}
(\Theta'_P(\rho_-,a),\Theta'_P(\rho_-,b))
\simeq
\begin{cases}
\mathrm{R}\Gamma(B,\mathcal L_{a-b}), & a\le b,\\
0, & a>b.
\end{cases}
\]
Here the inequality is reversed because
\[
(-\infty,a]\subset(-\infty,b]
\]
if and only if \(a\le b\).

This example shows the main point of the toric fibration version of the CCC.
The fan and the FLTZ skeleton still control the microlocal directions. The
theta sheaves are still supported on the same polyhedra in \(M_{\mathbb R}\).
The Hom-complex is still governed by inclusion of these polyhedra. What changes
is the coefficient of a nonzero Hom-complex: instead of obtaining \(k\), one
obtains the derived global sections of an associated line bundle on the base.

In this sense, the toric fibration version is a relative form of the ordinary
toric CCC:
\[
\text{polyhedral inclusion}
\quad\Longleftrightarrow\quad
\text{regular monomial direction},
\]
but the monomial direction now carries the additional coefficient
\(\mathcal L_\lambda\) coming from the principal torus bundle over \(B\).
When \(B=\operatorname{Spec}\mathbb C\), all associated line bundles have
one-dimensional spaces of global sections, and the ordinary toric CCC is
recovered.



\subsection{Fourier--Mukai Transforms from the Constructible Viewpoint}

The third generalization concerns Fourier--Mukai transforms. In the coherent
setting, a Fourier--Mukai transform is a way to describe a functor between
derived categories by using a kernel on a product. This should be understood
as a categorical form of correspondence. Instead of studying only ordinary maps
\(X\to Y\), one studies objects on \(X\times Y\), which act as generalized
relations between \(X\) and \(Y\).

Let \(X\) and \(Y\) be smooth projective varieties, and let
\[
K\in \mathrm{Perf}(X\times Y).
\]
We regard \(K\) as a kernel from \(X\) to \(Y\). It defines a functor
\[
\Phi_K:\mathrm{Perf}(X)\longrightarrow \mathrm{Perf}(Y)
\]
by
\[
\Phi_K(E)
=
\mathrm{R}p_{Y*}\bigl(p_X^*E\otimes^{\mathrm L}K\bigr),
\]
where \(p_X:X\times Y\to X\) and \(p_Y:X\times Y\to Y\) are the two projections.
Thus the construction has three steps: pull \(E\) back to \(X\times Y\), tensor
it with the kernel \(K\), and then push the result forward to \(Y\).

The reason that \(K\) may be viewed as a correspondence is already visible for
ordinary morphisms. If \(f:X\to Y\) is a morphism, its graph
\[
\Gamma_f\subset X\times Y
\]
is a subvariety of the product. The structure sheaf \(\mathcal O_{\Gamma_f}\)
is a kernel, and the associated Fourier--Mukai transform recovers the derived
pushforward:
\[
\Phi_{\mathcal O_{\Gamma_f}}(E)
\simeq
\mathrm{R}f_*E.
\]
Thus an ordinary map gives a special kind of kernel. Fourier--Mukai transforms
generalize this by allowing kernels that do not necessarily come from graphs.

Several standard functors are of Fourier--Mukai type. We record the basic
examples that will be useful for comparison.

\begin{enumerate}
    \item The identity functor is given by the diagonal kernel. If
    \(\Delta:X\hookrightarrow X\times X\) is the diagonal embedding, then
    \[
    K=\Delta_*\mathcal O_X
    \]
    defines
    \[
    \Phi_K\simeq \mathrm{id}_{\mathrm{Perf}(X)}.
    \]
    Thus the diagonal plays the role of the identity correspondence.

    \item Tensoring by a line bundle is also a Fourier--Mukai transform. If
    \(L\) is a line bundle on \(X\), then the kernel
    \[
    K=\Delta_*L\in \mathrm{Perf}(X\times X)
    \]
    defines the functor
    \[
    E\longmapsto E\otimes L.
    \]
    This is the example that we will compute below in the toric setting.

    \item Pushforward along a morphism is obtained from the graph kernel. Let
    \(f:X\to Y\) be a morphism, and let
    \[
    \Gamma_f\subset X\times Y
    \]
    be its graph. Then
    \[
    K=\mathcal O_{\Gamma_f}\in \mathrm{Perf}(X\times Y)
    \]
    defines
    \[
    \Phi_K(E)\simeq \mathrm{R}f_*E.
    \]
    Hence an ordinary map \(f:X\to Y\) is a special case of a
    Fourier--Mukai correspondence.

    \item Pullback along a morphism is obtained from the same graph, but with
    the two factors reversed. If \(f:X\to Y\), then we regard
    \(\Gamma_f\) as a subvariety of \(Y\times X\). The corresponding kernel
    gives a functor
    \[
    \mathrm{Perf}(Y)\longrightarrow \mathrm{Perf}(X),
    \]
    which is identified, under the usual derived convention, with
    \[
    F\longmapsto \mathrm{L}f^*F.
    \]
    Thus pullback is also represented by a kernel, provided the derived
    pullback is defined in the chosen category.

    \item More generally, a geometric correspondence gives a Fourier--Mukai
    type functor. Suppose that
    \[
    X \xleftarrow{p} Z \xrightarrow{q} Y
    \]
    is a correspondence and that \(F\in \mathrm{Perf}(Z)\). If
    \((p,q):Z\to X\times Y\), then the kernel
    \[
    K=\mathrm{R}(p,q)_*F
    \]
    defines a functor of the form
    \[
    E\longmapsto \mathrm{R}q_*(p^*E\otimes^{\mathrm L}F),
    \]
    whenever these operations are well-defined. Thus a Fourier--Mukai kernel
    can be viewed as a categorical correspondence between \(X\) and \(Y\).
\end{enumerate}

The word ``Fourier'' is most visible in the classical case of abelian varieties:
there the kernel is the Poincaré line bundle on \(A\times \widehat A\), where
\(\widehat A\) is the dual abelian variety. This kernel plays the role of the
exponential kernel in the classical Fourier transform. In modern usage,
however, the term Fourier--Mukai transform is used more generally for functors
defined by kernels.

From the point of view of the coherent--constructible correspondence, the
natural question is how such kernel functors appear on the constructible side.
Since the CCC identifies coherent theta objects with constructible theta
sheaves,
\[
\Theta'(\sigma,\chi)\longleftrightarrow \Theta(\sigma,\chi),
\]
one can study a Fourier--Mukai type functor by computing its action on theta
generators. If the coherent kernel sends theta objects to explicitly described
coherent objects, then the CCC translates this action into an operation on
polyhedral sheaves.

The example below is different from the toric fibration example of the previous
subsection. There, the main point was that nonzero Hom-complexes acquire
coefficients from the base. Here we keep a single toric variety and study how a
kernel acts as a functor on the theta generators.

Let \(X=\mathbb P^2\). Its fan has rays generated by
\[
v_1=(1,0),\qquad
v_2=(0,1),\qquad
v_3=(-1,-1),
\]
with maximal cones
\[
\sigma_{12}=\langle v_1,v_2\rangle,\qquad
\sigma_{23}=\langle v_2,v_3\rangle,\qquad
\sigma_{31}=\langle v_3,v_1\rangle.
\]
Consider the Fourier--Mukai transform given by tensoring with the equivariant
line bundle
\[
\mathcal O_{\mathbb P^2}(1)\cong \mathcal O_{\mathbb P^2}(D_{\rho_3}).
\]
Equivalently, the kernel is
\[
K=\Delta_*\mathcal O_{\mathbb P^2}(1)
\in \mathrm{Perf}(\mathbb P^2\times\mathbb P^2).
\]

To compute its effect on theta generators, we use the local Cartier data of
\(\mathcal O_{\mathbb P^2}(1)\). With the above convention, these local
characters are
\[
m_{\sigma_{12}}=(0,0),\qquad
m_{\sigma_{23}}=(1,0),\qquad
m_{\sigma_{31}}=(0,1).
\]
Indeed, for \(D_{\rho_3}\), the defining condition is
\[
\langle m_\sigma,v_\rho\rangle=-c_\rho
\]
for rays \(\rho\subset\sigma\), where \(c_{\rho_3}=1\) and
\(c_{\rho_1}=c_{\rho_2}=0\).

On the coherent side, tensoring by \(\mathcal O_{\mathbb P^2}(1)\) shifts the
local equivariant weight on each affine chart:
\[
\Theta'(\sigma,\chi)
\longmapsto
\Theta'(\sigma,\chi+m_\sigma).
\]
Under the CCC, this becomes
\[
\Theta(\sigma,\chi)
\longmapsto
\Theta(\sigma,\chi+m_\sigma).
\]
Since
\[
\Theta(\sigma,\chi)
=
(i_{\chi+\sigma^\vee})_*\omega_{\chi+\sigma^\vee},
\]
the corresponding constructible operation is the translation
\[
\chi+\sigma^\vee
\longmapsto
\chi+m_\sigma+\sigma^\vee.
\]

Let us spell this out for the three maximal cones.

For \(\sigma_{12}\), one has
\[
\sigma_{12}^\vee
=
\{m_1\ge 0,\ m_2\ge 0\}.
\]
Thus
\[
P_{\sigma_{12},(a,b)}
=
\{m_1\ge a,\ m_2\ge b\}.
\]
Since \(m_{\sigma_{12}}=(0,0)\), tensoring by
\(\mathcal O_{\mathbb P^2}(1)\) does not shift the theta sheaves on this chart:
\[
P_{\sigma_{12},(a,b)}
\longmapsto
P_{\sigma_{12},(a,b)}.
\]

For \(\sigma_{23}\), the dual cone is
\[
\sigma_{23}^\vee
=
\{m_2\ge 0,\ m_1+m_2\le 0\}.
\]
Hence
\[
P_{\sigma_{23},(a,b)}
=
\{m_2\ge b,\ m_1+m_2\le a+b\}.
\]
Since \(m_{\sigma_{23}}=(1,0)\), the transform sends this sector to
\[
P_{\sigma_{23},(a+1,b)}
=
\{m_2\ge b,\ m_1+m_2\le a+b+1\}.
\]
Thus the corresponding constructible theta sheaf is translated by the local
Cartier character \((1,0)\).

For \(\sigma_{31}\), the dual cone is
\[
\sigma_{31}^\vee
=
\{m_1\ge 0,\ m_1+m_2\le 0\}.
\]
Hence
\[
P_{\sigma_{31},(a,b)}
=
\{m_1\ge a,\ m_1+m_2\le a+b\}.
\]
Since \(m_{\sigma_{31}}=(0,1)\), the transform sends this sector to
\[
P_{\sigma_{31},(a,b+1)}
=
\{m_1\ge a,\ m_1+m_2\le a+b+1\}.
\]

This calculation shows how the diagonal kernel
\[
\Delta_*\mathcal O_{\mathbb P^2}(1)
\]
appears on the constructible side. It acts on coherent theta generators by
shifting their local equivariant weights. Under the CCC, the same operation
becomes translation of the polyhedral supports of the corresponding
constructible theta sheaves.

The Hom-complex pattern is compatible with this operation. If
\[
P_{\sigma,\chi}\subset P_{\sigma,\psi},
\]
then after tensoring by \(\mathcal O_{\mathbb P^2}(1)\) the corresponding
inclusion becomes
\[
P_{\sigma,\chi+m_\sigma}\subset P_{\sigma,\psi+m_\sigma}.
\]
Thus the inclusion condition is preserved. Equivalently, on the coherent side,
tensoring by a line bundle is an autoequivalence, and the same autoequivalence
is seen constructibly as translation of the theta polyhedra.

Therefore this subsection illustrates a different use of the CCC from the
toric fibration case. A toric fibration changes the values of the Hom-complexes
by introducing coefficients from a base. A Fourier--Mukai kernel, by contrast,
acts as a functor on the CCC category itself. In the example above, the kernel
\(\Delta_*\mathcal O_{\mathbb P^2}(1)\) acts by translating the polyhedral
supports determined by the local Cartier data of the line bundle.
